\chapter{Conclusion and Future Scope}
\label{ch:conclusion}

\section{Conclusion}

Pure Transformers become slow and memory-heavy on long inputs because
self-attention compares every pair of tokens and keeps a KV cache that grows
with the context. State-space models such as Mamba-2 avoid both costs by
carrying a constant-size state, but that state is a lossy summary. This
seminar studied \NS, a model that combines the two: 40 Mamba-2 blocks carry
most of the sequence processing, 8 attention layers act as global anchors for
exact long-range retrieval, and 40 LatentMoE layers provide capacity.

Three further ideas complete the design. LatentMoE moves the experts into a
1024-dimensional latent space, which cuts memory and routing traffic fourfold
and pays for more experts, so that 120 billion parameters exist while only
about 12 billion are active for each token. Shared-weight multi-token
prediction heads give the model its own draft mechanism for speculative
decoding. Pretraining in 4-bit NVFP4 shows that low-precision training is
stable even over 25 trillion tokens.

The published results show a model in the same accuracy class as the best
open Transformer models of its size---ahead of GPT-OSS-120B on most tasks and
somewhat behind Qwen3.5-122B---while running up to 2.2$\times$ and
7.5$\times$ faster than them and holding a RULER score of 91.6 at one million
tokens. The study also showed that the hybrid approach is not free: it brings
new engineering problems, of which the accumulation of rounding error in the
Mamba state is the clearest example.

The broader lesson is that attention does not have to be everywhere. Used
sparingly and placed carefully, a handful of attention layers appear to be
enough, and the rest of the network can be built from cheaper components.
\NS\ is strong evidence that the next generation of efficient large-scale
models will be hybrids designed together with the hardware they run on, rather
than pure Transformers.

\section{Future Scope}

Several directions follow from this study.
\begin{itemize}
  \item \textbf{Larger hybrids.} The Nemotron~3 family includes a larger
        \emph{Ultra} model~\cite{nvidia2025nemotron3}; it will show whether the
        same layout continues to scale.
  \item \textbf{Placement and number of attention layers.} The ratio of about
        8--9\% attention comes from experiments on small models. A systematic
        study at large scale, or a learned placement, could improve the
        trade-off further.
  \item \textbf{Quantization-robust state-space layers.} The Mamba state
        currently needs a training-free fix at deployment. Training methods
        that make the recurrence itself robust to low-precision storage would
        remove a source of fragility and might allow the state to be stored in
        fewer bits.
  \item \textbf{Lower-precision training everywhere.} Several layer types are
        still trained in BF16. Extending 4-bit training to them, and
        understanding the growth of zero-valued gradients, would reduce
        training cost further.
  \item \textbf{Independent benchmarking.} Reproducing the throughput results
        on other accelerators and serving stacks would clarify how much of the
        advantage is architectural and how much is specific to one hardware
        platform.
  \item \textbf{Other modalities.} Linear-time sequence layers are attractive
        wherever sequences are very long, such as audio, video and genomic
        data; hybrid backbones could be extended to multimodal models.
  \item \textbf{Smaller footprints.} Techniques for loading experts on demand
        or pruning rarely used experts could reduce the memory needed to hold
        all 120 billion parameters and bring such models to smaller
        deployments.
\end{itemize}
